\documentclass[main]{subfiles}

\title[Advanced Hyperparameter Optimization]{Advanced Hyperparameter Optimization}
\subtitle{Warmstarting}
\author[Katharina Eggensperger]{Katharina Eggensperger}
\date{}

\titlebackground*{images/AutoML_Speeding_up_HPO.jpg}

%\institute{}
%\date{}

\begin{document}
	
	\maketitle
	

%----------------------------------------------------------------------
%----------------------------------------------------------------------
\begin{frame}[c]{Why Should We Start From Scratch?}

In practice, we have more data available from:
\begin{itemize}
    \item tune the same model repeatedly
    \item previous work on similar datasets
    \item re-run experiments with slight modifications
\end{itemize}

\vspace{1cm}
\textbf{$\leadsto$ How can we use this information to increase the efficiency of Bayesian optimization?} \\

\end{frame}
%-----------------------------------------------------------------------
%----------------------------------------------------------------------
\begin{frame}[c]{What Is Warmstarting?}

In Bayesian optimization, we iteratively build a surrogate model: 

\begin{equation*}
    \surro_n (\conf | \hist), \hist = D_{init}0 \cup \{(\conf_i, \cost_i)\}_{i=n}^N
\end{equation*}

\begin{itemize}
    \item efficiency and quality of early iterations are highly sensitive to the initial design $\dataset_{init}$.
    \item standard practice: Random or space-filling designs (e.g., Latin Hypercube)
\end{itemize}

\vspace{1cm}
\textbf{$\leadsto$ How do we replace uninformed initialization with informed priors over good regions of $\pcs$?}

\end{frame}
%-----------------------------------------------------------------------%----------------------------------------------------------------------
\begin{frame}[c]{Let's Consider Two Settings}

Assuming a fixed searchspace and given a target dataset $\dataset_t$, we consider two settings

\textbf{Simple Case:} We have prior observations on the \textit{same} dataset (e.g., from a manual exploration)

\textbf{Advanced Case:} We have prior observations on \textit{related} datasets $\{\dataset_1, \dots, \dataset_{t-1}\}$ (e.g., the same HPO task, but on a different dataset)

\vspace{0.5cm}
$\leadsto$ inject knowledge into the selection of the next configuration to evaluate in the early stages of Bayesian Optimization.

\setbeamercolor{block title}{use=normal text,fg=normal text.fg}%
\setbeamercolor{block body}{use=normal text,fg=normal text.fg}%
\begin{block}{\raggedright Impact}
  \begin{itemize}
    \item reduces initial uncertainty volume
    \item biases surrogate towards existing (good) observations
    \item accelerates early exploitation
  \end{itemize}
\end{block}

\end{frame}
%-----------------------------------------------------------------------
%----------------------------------------------------------------------
\begin{frame}{Simple Case: Prior Observations on the Same Task}

\textbf{Assumption:} Prior observations are based on human intuition and serve as good starting points. \\
\textbf{Solution:} (Partially) replace initial design iterations with observations from manual tuning or prior experiments:

{%
\setbeamercolor{block title}{use=normal text,fg=normal text.fg}%
\setbeamercolor{block body}{use=normal text,fg=normal text.fg}%
\begin{block}{\raggedright Limitation}
  \begin{itemize}
    \item does not work across tasks
    \item performance depends on prior observations
  \end{itemize}
\end{block}
}%

\end{frame}
%-----------------------------------------------------------------------
%----------------------------------------------------------------------
\begin{frame}{Advanced Case: Prior Observations on Related Tasks}

\textbf{Assumption:} Objective functions of tasks are correlated. \\
\textbf{Challenge:} Positive transfer from similar tasks, zero transfer from different tasks. \\
\textbf{Solution:} Portfolios

\vspace{1cm}

\textbf{$\leadsto$Direction 1:} \textbf{Task-agnostic} portfolios covering well-performing regions of "all" tasks. \\
\textbf{$\leadsto$Direction 2:} \textbf{Task-specific} portfolios covering well-performing regions on similar tasks.

\end{frame}
%-----------------------------------------------------------------------
%----------------------------------------------------------------------
\begin{frame}{Task-Agnostic Portfolios}

We assume access to datasets $\datasets_{meta} = \dataset_{1:t-1}$ and a candidate pool $\confs_{cand} = \{\conf_1, ..., \conf_n\}$. The goal is to construct a portfolio $\portfolio \subseteq \confs_{cand}, |\portfolio| = k$, that minimizes the objective:

\begin{equation*}
\loss(\portfolio, \datasets_{meta}) = \sum_{i=1}^{t-1} \min_{\conf \in \portfolio} \cost(\conf, \dataset_i)
\end{equation*}

\textbf{Key idea:}
\begin{itemize}
    \item Portfolio performance is determined by the single best configuration.
    \item The portfolio should include at least one well-performing configuration per task.
    \item The portfolio covers optima across diverse tasks.
\end{itemize}

\end{frame}
%-----------------------------------------------------------------------
%----------------------------------------------------------------------
\begin{frame}{Greedy Forward Selection}

\textbf{Intuition:}
\begin{itemize}
    \item Each step adds a configuration that improves \emph{uncovered tasks}
    \item Encourages diversity via complementary performance
\end{itemize}

\vspace{0.5cm}

\begin{columns}[T,onlytextwidth]
\begin{column}{0.48\textwidth}
\textbf{Input:} $\confs_{cand} = \{\conf_1, \dots, \conf_n\}$, $\dataset_{1:t-1}$, $\cost(\cdot)$ \\
\textbf{Initialize:} $\portfolio \gets \emptyset$ \\
\begin{enumerate}
    \item \textbf{while $|\portfolio| < k|$ do} \\
    \item \quad $\conf^+ \in \argmin_{\lambda \in \confs_{cand}} \loss \left(\portfolio \cup \conf, \datasets_{meta} \right)$ \\
    \item \quad $\portfolio \gets \portfolio \cup \conf^+, \confs_{cand} \gets \confs_{cand} \setminus \conf^+$ \\
    \item \textbf{return $\portfolio$}
\end{enumerate}
\end{column}
\begin{column}{0.48\textwidth}
{%
\setbeamercolor{block title}{use=normal text,fg=normal text.fg}%
\setbeamercolor{block body}{use=normal text,fg=normal text.fg}%
\begin{block}{\raggedright Limitation}
  \begin{itemize}
    \item requires a set of candidate configurations
    \item required performance of all configurations on all tasks
    \item not all tasks are relevant for the current task
  \end{itemize}
\end{block}
}%
\end{column}
\end{columns}
\vspace{0.5cm}

\end{frame}
%-----------------------------------------------------------------------
%----------------------------------------------------------------------
\begin{frame}{Task-Specific Portfolios}

Instead of finding a global portfolio, we want a portfolio specifically for our target task $t$.

\vspace{0.5cm}
\textbf{Key idea} on how to adapt greedy portfolio construction:
\begin{itemize}
    \item Reduce the set of meta-datasets to only similar tasks.
    \item Select the best performing configurations on these tasks as candidate configurations.
    \item Construct portfolio.
\end{itemize}
\vspace{0.5cm}

\textbf{How do we define task similarity?} \\
\pause
$\leadsto$ Intuition: If tasks are similar the similar configurations achieve similar performance, i.e., if the datasets are similar, good configurations transfer.

\end{frame}
%-----------------------------------------------------------------------
%----------------------------------------------------------------------
\begin{frame}{Meta-features as Dataset Descriptors}

We assume access to dataset descriptors (meta-features) $\mathbf{z} = \{\phi_1(\dataset) \dots, \phi_d(\dataset)\}$ representing dataset characteristics, e.g.,
\begin{itemize}
    \item basic features (e.g., \#samples, \#features),
    \item information-theoretic meta-features (e.g., entropy of features and labels),
    \item landmarking features (e.g., performance of simple, fast ML models),
    \item and learned meta-features (e.g., task/dataset embeddings).
\end{itemize}
Task similarity can then be computed as $\text{sim}(t, t') = -\| \mathbf{z}_t - \mathbf{z}_{t'}\|$
\end{frame}
%-----------------------------------------------------------------------
%----------------------------------------------------------------------
\begin{frame}{Meta-feature-based Portfolios}

We can now build task-specific portfolios:
\begin{enumerate}
    \item Compute $\{\mathbf{z}_1, \cdots, \mathbf{z}_t\}$
    \item Use the best performing configurations of the closest $k$ datasets as the portfolio
\end{enumerate}

\setbeamercolor{block title}{use=normal text,fg=normal text.fg}%
\setbeamercolor{block body}{use=normal text,fg=normal text.fg}%
\begin{block}{\raggedright Limitation}
\begin{itemize}
    \item Requires hand-designed meta-features
    \item Similarity may not align with the optimization landscape
\end{itemize}
\end{block}
   
\end{frame}
%----------------------------------------------------------------------
%----------------------------------------------------------------------
\begin{frame}{Alternatives}

Alternative directions to inject knowledge observed on related tasks:
\begin{itemize}
    \item Transfer hyperparameters (e.g., lengthscales) from related tasks.
    \item Combine/ensemble surrogate models build on related tasks.
    \item Measure similarity by comparing the ranking of configurations.
\end{itemize}

\end{frame}
%----------------------------------------------------------------------
%----------------------------------------------------------------------
\begin{frame}{Conclusion}

\textbf{Summary:}
\begin{itemize}
    \item Warmstarting Bayesian optimization increases the efficiency of HPO.
    \item We can (partially) replace the initial design with a portfolio.
    \item Portfolios can be task-dependend or task-agnostic.
\end{itemize}

\textbf{Takeaway:}
\begin{itemize}
    \item Better warmstarting requires stronger assumptions about task similarity. 
    \item Stronger assumptions increase inductive bias. If that bias is correct, optimization becomes more efficient; when it is wrong, we get negative transfer.
    \item Warmstarting turns HPO into a meta-learning problem across datasets.
\end{itemize}

\end{frame}


\end{document}
